%! Transformations of Functions — Unit 1, Lesson 1.4 — Group Activity
\documentclass[10pt]{article}
\usepackage{precalculus-article}
\usepackage{precalculus-boxes}
\newcommand{\TallMath}[1]{$\displaystyle #1\rule[-1.4em]{0pt}{3.2em}$}

\begin{document}

\pageheader{Unit 1, Lesson 1.4}{Group Activity}

\vspace{0.12in}

\begin{scenariobox}[The Skate Park Ramp]{plumlight}
\small
A skate park's ramp has a curved cross-section modeled by the parent function
$f(x) = x^2$, where $x$ is the horizontal distance in feet from the low point
and $f(x)$ is the height of the ramp surface in feet. The design team keeps
asking for changes --- but the \textit{curve} never changes, only where it
sits and how big it is. Every tier uses this same table and pre-drawn graph.

\smallskip
\begin{center}
\renewcommand{\arraystretch}{1.4}
\begin{tabular}{|c|c|c|c|c|c|}
\hline
$x$ (feet) & $-2$ & $-1$ & 0 & 1 & 2 \\
\hline
$f(x)$ (feet) & 4 & 1 & 0 & 1 & 4 \\
\hline
\end{tabular}
\end{center}

\begin{center}
\begin{tikzpicture}[x=0.9cm, y=0.42cm]
  \draw[linegray, very thin, step=1] (-2,0) grid (2,5);
  \draw[->, thick] (-2.4,0) -- (2.5,0) node[below right] {\small $x$ (ft)};
  \draw[->, thick] (0,0) -- (0,5.5) node[above] {\small $f(x)$ (ft)};
  \foreach \x in {-2,-1,1,2} \node[below, font=\small] at (\x,0) {\x};
  \foreach \y in {1,2,3,4,5} \node[left, font=\small] at (0,\y) {\y};
  \draw[very thick, plum] plot [smooth] coordinates
    {(-2,4) (-1,1) (0,0) (1,1) (2,4)};
  \fill[plum] (-2,4) circle (2.5pt);
  \fill[plum] (-1,1) circle (2.5pt);
  \fill[plum] (0,0) circle (2.5pt);
  \fill[plum] (1,1) circle (2.5pt);
  \fill[plum] (2,4) circle (2.5pt);
\end{tikzpicture}
\end{center}
\end{scenariobox}

\vspace{0.1in}

\boxguard[30]
\begin{tcolorbox}[colback=white, colframe=black!40,
    title=\textbf{Tier R --- Remediate},
    fonttitle=\bfseries, arc=2mm, left=3mm, right=3mm, top=2mm, bottom=2mm]

\begin{enumerate}[itemsep=8pt]
  \item \textbf{Build it on a platform.} The crew raises the whole ramp onto a
        3-foot platform: $g(x) = f(x) + 3$. Complete the table.

        \smallskip
        \renewcommand{\arraystretch}{1.5}
        \begin{center}
        \begin{tabular}{|c|c|c|c|c|c|}
        \hline
        $x$ & $-2$ & $-1$ & 0 & 1 & 2 \\
        \hline
        $f(x)$ & 4 & 1 & 0 & 1 & 4 \\
        \hline
        $g(x) = f(x) + 3$ & \blank{1.1cm} & \blank{1.1cm} & \blank{1.1cm}
          & \blank{1.1cm} & \blank{1.1cm} \\
        \hline
        \end{tabular}
        \end{center}

        The whole ramp moved \textbf{up} / \textbf{down} (circle one) by
        \blank{1.2cm} feet.

  \item \textbf{Sink it in a pit.} Now the crew drops the ramp 1 foot into a
        pit: $h(x) = f(x) - 1$. Complete the table.

        \smallskip
        \renewcommand{\arraystretch}{1.5}
        \begin{center}
        \begin{tabular}{|c|c|c|c|c|c|}
        \hline
        $x$ & $-2$ & $-1$ & 0 & 1 & 2 \\
        \hline
        $h(x) = f(x) - 1$ & \blank{1.1cm} & \blank{1.1cm} & \blank{1.1cm}
          & \blank{1.1cm} & \blank{1.1cm} \\
        \hline
        \end{tabular}
        \end{center}

        The whole ramp moved \textbf{up} / \textbf{down} (circle one) by
        \blank{1.2cm} feet.

  \item \textbf{The sorting question.} In both rules, the ``$+3$'' and the
        ``$-1$'' were written \textbf{outside} / \textbf{inside} (circle one)
        the function $f$.

        So they changed the \blank{2cm} column of the table, and the ramp
        moved \blank{2.6cm}.
\end{enumerate}
\end{tcolorbox}

\vspace{0.1in}

\boxguard[30]
\begin{tcolorbox}[colback=white, colframe=black!40,
    title=\textbf{Tier A --- Approaching Proficiency},
    fonttitle=\bfseries, arc=2mm, left=3mm, right=3mm, top=2mm, bottom=2mm]

The design is revised to $g(x) = 2f(x - 1)$: the ramp is moved down the park
and made twice as steep.

\begin{enumerate}[itemsep=8pt]
  \item Complete the table \textbf{one step at a time} --- do the inside
        change first, then the outside one.

        \smallskip
        \renewcommand{\arraystretch}{1.5}
        \begin{center}
        \begin{tabular}{|c|c|c|c|c|c|}
        \hline
        $x$ & $-1$ & 0 & 1 & 2 & 3 \\
        \hline
        $x - 1$ & $-2$ & $-1$ & 0 & 1 & 2 \\
        \hline
        $f(x-1)$ & \blank{1.1cm} & \blank{1.1cm} & \blank{1.1cm}
          & \blank{1.1cm} & \blank{1.1cm} \\
        \hline
        $g(x) = 2f(x-1)$ & \blank{1.1cm} & \blank{1.1cm} & \blank{1.1cm}
          & \blank{1.1cm} & \blank{1.1cm} \\
        \hline
        \end{tabular}
        \end{center}

  \item Describe both moves in words --- name a direction and an amount for
        each.

        \vspace{0.05in}
        \writelines{2}

  \item The low point of the original ramp is at $(0, 0)$. The low point of
        the new ramp $g$ is at $($ \blank{1cm} , \blank{1cm} $)$.

  \item On the blueprint, $f$ has domain $[-2, 2]$ and range $[0, 4]$. Find
        the domain and range of $g$.

        \begin{work}
          \text{Domain: } -2 &\le x - 1 \le 2 \\
                          -1 &\le x \le 3 \\
          \text{Range: } 0 &\le y \le 4 \\
                         0 &\le 2y \le 8
        \end{work}

        Domain of $g$: \blank{2.6cm} \qquad Range of $g$: \blank{2.6cm}
\end{enumerate}
\end{tcolorbox}

\vspace{0.1in}

\boxguard[30]
\begin{tcolorbox}[colback=white, colframe=black!40,
    title=\textbf{Tier E --- Extension},
    fonttitle=\bfseries, arc=2mm, left=3mm, right=3mm, top=2mm, bottom=2mm]

\begin{enumerate}[itemsep=8pt]
  \item \textbf{Reverse-engineer it.} A second ramp $k$ is a transformation of
        $f$. Here are its values.

        \smallskip
        \renewcommand{\arraystretch}{1.4}
        \begin{center}
        \begin{tabular}{|c|c|c|c|c|c|}
        \hline
        $x$ & $-2$ & $-1$ & 0 & 1 & 2 \\
        \hline
        $k(x)$ & $-8$ & $-2$ & 0 & $-2$ & $-8$ \\
        \hline
        \end{tabular}
        \end{center}

        Write $k$ in terms of $f$: \quad $k(x) = $ \blank{3cm}

        Describe the two things that happened to the graph:

        \vspace{0.05in}
        \writeline

  \item A third design $m$ has these values.

        \smallskip
        \renewcommand{\arraystretch}{1.4}
        \begin{center}
        \begin{tabular}{|c|c|c|c|c|c|}
        \hline
        $x$ & 0 & 1 & 2 & 3 & 4 \\
        \hline
        $m(x)$ & 5 & 2 & 1 & 2 & 5 \\
        \hline
        \end{tabular}
        \end{center}

        The low point moved from $(0,0)$ to $($ \blank{1cm} , \blank{1cm} $)$,
        so \quad $m(x) = $ \blank{4cm}

  \item \textbf{Inside is not the same as outside.} Use $f(x) = x^2$ and the
        extra value $f(0.5) = 0.25$ to fill in this comparison.

        \smallskip
        \renewcommand{\arraystretch}{1.5}
        \begin{center}
        \begin{tabular}{|c|c|c|}
        \hline
        $x$ & $f(2x)$ & $2f(x)$ \\
        \hline
        $0.5$ & \blank{1.4cm} & \blank{1.4cm} \\
        \hline
        $1$ & \blank{1.4cm} & \blank{1.4cm} \\
        \hline
        \end{tabular}
        \end{center}

        Multiplying \textit{inside} and multiplying \textit{outside} give
        different graphs. Explain what each one does.

        \vspace{0.05in}
        \writelines{2}
\end{enumerate}
\end{tcolorbox}

\end{document}
